\documentclass{amsart}

\usepackage{amsmath,amsthm,amssymb}
\usepackage{hyperref}
\usepackage{paralist}

\usepackage{enumitem}
%\usepackage[pagewise]{lineno}\linenumbers
\usepackage[pagewise]{lineno}
\usepackage[title]{appendix}
\usepackage[framemethod=TikZ]{mdframed}
\usepackage[labelformat=empty]{subfig}
\usepackage[formats]{listingsutf8}

\makeatletter
\DeclareFontFamily{U}{tipa}{}
\DeclareFontShape{U}{tipa}{bx}{n}{<->tipabx10}{}

\usepackage{mdframed}

\newcommand{\abs}[1]{\left|#1\right|}  %\abs{x} yields |x|.
\newcommand{\assign}{\mathrel{\mathop :}=}
\DeclareMathOperator{\sgn}{sgn}
\DeclareMathOperator{\rms}{rms}

\setlist{itemsep=4pt}
\renewcommand{\lstlistingname}{Script}
\lstset{
    literate = {Ω}{{$\Omega$}}1,
    language = Matlab,
    basicstyle = \scriptsize,
    tabsize = 4,
    frame = single,
    framerule = 0.5pt,
    captionpos = b,
    keywordstyle = \ttfamily\color{magenta},
    identifierstyle = \ttfamily\color{cyan},
    commentstyle = \ttfamily\color{gray},
    stringstyle = \ttfamily\color{violet},
    showstringspaces = false,
    backgroundcolor = \color{yellow!10},
    breaklines = true
}

% Set up problem environment.
\newtheorem{theorem}{Theorem}
\newtheorem{Problem}{Problem}

\usepackage[title]{appendix}

\begin{document}

\title[aps1]{Lab Report 1 [{\bf lab1}]\\
insert Your name\\
insert Your Student Number\\
insert Your email address\\
insert Current date and time}

%\author{Your Name}
%\address{
%Student number\\
%Your email address
%}
%\email{you@umanitoba.ca}

\date{2023-09-15}

\begin{abstract}
In preparing your solutions to Lab 1 (\colorbox{green!15}{\bf lab1}), type your solutions using either MS Word (and MathType for formulas) or \LaTeX\ (it does both text \& formulas \colorbox{red!15}{\bf and is used by Matlab} for concise rendition of formulas in labelling plot axes, legends and titles).  A template for your solutions written in \LaTeX\  accompanies the aps1 files.  All solutions must be typed.   In submitting your solutions for {\bf lab1}, use the UMLearn dropbox for lab1.
The deadline for {\bf lab1} is [\colorbox{magenta!15}{\bf 3 Feb. 2025}].  
\end{abstract}

\keywords{Euler Exponential, Frequency (via Mersenne Primes), Matlab, Pixel, Polar Form, Scalar Field, Digital Image, Vector}
\maketitle

%\pagebreak

\tableofcontents
\thispagestyle{empty}

\section{Instructions}
In your Assessment Problems Set (aps) solutions report, do the following:
\begin{compactenum}[{Do}1$^o$]
\item Always start your aps1 solutions report with your
\begin{compactenum}[$\bullet$.1]
\item Your full name.
\item Your Student Number.
\item Your email address.
\item Current date and time.
\end{compactenum}
\item Using either MS Word (and MathType for formulas) or \LaTeX, type each of your answers.
\colorbox{magenta!15}{Handwritten answers will not be accepted.}
\item Use either a drawing tool or Matlab to produce plots for a problem solution.
\item Insert an Appendix [Matlab Scripts] in your aps1 report that gives a listing of each Matlab script used in your report.  
\item \colorbox{brown!15}{\bf State the problem} in the aps, before you give your solution.
\item Upload your aps solutions report to the UMLearn dropbox at \url{https://universityofmanitoba.desire2learn.com/d2l/home/617957}.  Your report should include your (1) problem solutions and (2) all Matlab scripts you used to obtain your solutions.
\item \colorbox{green!15}{Matlab script naming convention}.  In naming each Matlab script, insert the name of the problem, e.g., for the solution of Problem 3, use \colorbox{brown!15}{lab1-3} as the name of your script.
\item At the end of your lab1 report, insert
\begin{compactenum}[(i)]
\item short biography about yourself.
\item Picture of yourself.
\end{compactenum}
\end{compactenum}

\section{Problems}

\begin{compactenum}[\colorbox{magenta!15}{\bf lab1-}1]
 \item Modify the exmple1LinewidthProb1.m Matlab script from Lecture 1 to obtain the required plot.

% ---------------------------------------------------
    
%%\item\label{ProbB22} 
\item Modify the exmpleExTanhProb2.m Matlab script from Lecture 1 to obtain the required plot.
\item Modify exmpleEulerProb3.m Matlab script from Lecture 1 to obtain the required plot.
\item Modify the testImgScalarField.m Matlab script from Lecture 1 to obtain the 3 required suplots.
\item Prove:\\
	\begin{theorem} Every digital image is a scalar field.
	\end{theorem}.
\item Prove:\\
	\begin{theorem} Every pixel in a digital image is a vector.
	\end{theorem}
\item In an Observations section of your report,
indicate anything special or interesting that you found in solving one or more of the lab problems.
\end{compactenum}

\begin{appendix}
\section{Sample Matlab script}
\lstinputlisting[]{./scripts/image3Types.m}

\end{appendix}

%\bibliographystyle{amsplain}
%\bibliography{NSrefs}

\end{document}
